% Standalone solution dossier for prop:a5-ratio.
\documentclass[../main.tex]{subfiles}
\begin{document}

% === SOLUTION HEADER (proof provenance) =====================================
% ledger-nodes: prop:a5-ratio; obs:symmetry-vs-physical
% refines: prop:a5-ratio
% derived: obs:symmetry-vs-physical
% bounded_by: none
% evidence_run: none
% checked_by: agent
% author: /root/a4_a5
% base reviewer: /root/review_a4_a5
% dossier_editor: /root/review_a4_a5
% base review: research/reviews/2026-08-21-a4-a5-independent-agent-audit.md
% derived reviewer: /root/audit_metadata
% derived review: research/reviews/2026-08-22-derived-obstructions-audit.md
% derived date: 2026-08-22
% date: 2026-08-21
% ============================================================================

\section*{Solution: invariant and non-invariant spectral blocks}

\begin{proposition}
Let a finite group $G$ preserve a probability measure $\pi$ and its Dirichlet form, and let
$\bar\pi$ carry the lifted quotient form. Define $P_Gf=|G|^{-1}\sum_g f\circ g$ and the
restricted gaps
\[
 \lambda_{\mathrm{inv}}=\inf_{\substack{f\in H^1(\pi)^G\\\Var_\pi(f)>0}}
 \frac{\mathcal E_\pi(f)}{\Var_\pi(f)},\qquad
 \lambda_{\mathrm{noninv}}=\inf_{\substack{f\in H^1(\pi),\ P_Gf=0\\\Var_\pi(f)>0}}
 \frac{\mathcal E_\pi(f)}{\Var_\pi(f)}.
\]
Then
\[
 C_P(\bar\pi)=\lambda_{\mathrm{inv}}^{-1},\qquad
 C_P(\pi)=\min\{\lambda_{\mathrm{inv}},\lambda_{\mathrm{noninv}}\}^{-1},
\]
with extended reciprocals. If $0<\lambda_{\mathrm{inv}}<\infty$, strict quotient improvement
holds exactly when $\lambda_{\mathrm{noninv}}<\lambda_{\mathrm{inv}}$.
\end{proposition}

\begin{proof}
Invariance makes every group action unitary for both the $L^2(\pi)$ and energy pairings. Hence
$P_G$ is the orthogonal projection in both pairings. Every mean-zero $f$ therefore decomposes as
$f=P_Gf+(I-P_G)f$ with
\[
 \Var_\pi(f)=\Var_\pi(P_Gf)+\Var_\pi((I-P_G)f),\qquad
 \mathcal E_\pi(f)=\mathcal E_\pi(P_Gf)+\mathcal E_\pi((I-P_G)f).
\]
Its Rayleigh quotient is a variance-weighted average of the two block quotients, so the full gap
is their minimum. Pullback by the quotient map is an isometry from $L^2(\bar\pi)$ onto the
invariant block and preserves the Dirichlet form, variance, and constants. Thus the quotient gap
is $\lambda_{\mathrm{inv}}$. Inversion gives the identities and, under the stated finite positive
condition, the strict-improvement criterion.
\end{proof}

\paragraph{Obstruction respected.}
A tied non-invariant eigenfunction does not imply improvement; the non-invariant gap must be
strictly smaller than the invariant one.

\end{document}
